\subsection*{\S. 13.\quad 管状面。}
作为第一个例子，考察一个半径为常数 \(a\) 的球面；它在空间中运动，形成一个单参数的位置族。我们来求它的包络面。为了定义这个 \(\infty^1\) 球面族，只需给出球心的轨迹曲线 C；对这条曲线仍采用通常的记号。动球面的方程为
\begin{equation}\tag{19}
(X-x)^2+(Y-y)^2+(Z-z)^2=a^2
\end{equation}
此时的变参数是曲线 C 的弧长 \(s\)。为了求球面 (19) 上的特征线，应把这个方程与它对 \(s\) 作一次微分所得的方程联立，即
\begin{equation}\tag{20}
(X-x)\alpha+(Y-y)\beta+(Z-z)\gamma=0,
\end{equation}
这就是曲线 C 的法平面方程。因此：\emph{动球面的特征线，是它的一个大圆；这个大圆位于曲线 C 在球心处的法平面内。}这个半径恒为 \(a\) 的动圆，其圆心描绘曲线 C，而其平面始终保持与该曲线垂直；因此，这个动圆的轨迹就是所求包络面。这样的曲面称为\emph{管状面}。

现在若要确定包络管状面上的回归棱，就必须把 (20) 再微分一次，并将这样所得的新方程
\begin{equation}\tag{21}
(X-x)\xi+(Y-y)\eta+(Z-z)\zeta=\rho
\end{equation}
与 (19)、(20) 联立。方程 (21) 正是这样一个平面的方程：它垂直于主法线，并经过主法线上与 \(\mathrm{M}\equiv(x,y,z)\) 相距恰为第一曲率半径 \(\rho\)、且位于正向一侧的点（曲率中心）。所以，当 \(\rho<a\) 时，特征圆上的极限点为相异的实点；当 \(\rho=a\) 时，它们为重合的实点；当 \(\rho>a\) 时，它们为虚点。

因此，在曲线 C 上满足 \(a>\rho\) 的区域内，管状面上的回归棱是实的，并由两个不同的分支组成。
